\section{The Flow}

The base is feeling, and only feeling. Everything in the paper is a pattern made of vibes, from a single tone up to a whole mind, and nothing is ever added underneath.

\begin{axiom}[The vibe]
The one substance is the vibe, which is experience. A vibe's tone is the quality of that experience, not a thing apart from it. There is nothing beneath it.
\end{axiom}

\begin{axiom}[The flow]
The whole is the flow, the eight elements running as one, written $F$, the vibe in endless motion.
\end{axiom}

The flow is the highest view of everything, the ceaseless changing the universe passes through from one moment to the next, the experiential field in active and endless motion, never still and never finished. The eight elements set down below are what it is made of, the pieces whose running together is the flow, and the definitions that follow name each of them in turn.

A map, though, is not the territory. Those eight definitions, and the whole of this paper, are a model of the structure of the flow and not the flow itself. The flow is experience in motion, and experience is not had by reading a description of it. It is had by living it. What the model gives is the shape of the thing, exact enough to compute and to test. The thing itself, the felt flow, can only be known directly, from the inside, in the raw experience of your own life.

To say what a vibe is and how it moves takes exactly eight words. Each is a short four-letter name, and each begins with its own letter, so it doubles as a clean symbol. Three give the where, a space and its parts. Two give the what, a value and its order. Two give the how, one law for the state and one for the space. The last gives the when.

\begin{definition}[Mesh]
The whole space, the growing $\{3,4,3,4\}$ hyperbolic honeycomb.
\end{definition}

\begin{definition}[Dock]
One cell of the mesh, a here-now where tones meet and move on.
\end{definition}

\begin{definition}[Site]
One of a dock's twenty-four directions, a $D_4$ root, where a single tone lives.
\end{definition}

\begin{definition}[Tone]
One way to model the quality of a vibe's experience, a value in $\{-1, 0, +1\}$, read as pain, peace, or pleasure. The tone is not separate from the vibe, and not a thing the vibe holds. It is that quality, named as a value.
\end{definition}

\begin{definition}[Lean]
The order $-1 < 0 < +1$, which turns a tone into signed charge and gives creation a direction.
\end{definition}

\begin{definition}[Knit]
The law of the state, collide then stream, once a beat.
\end{definition}

\begin{definition}[Wake]
The law of the space, new docks born at peace at the growing edge.
\end{definition}

\begin{definition}[Beat]
One discrete tick of time.
\end{definition}

The eight are not a loose pile. They obey hard constraints, and the constraints are what keep the model from cheating. A base element must be discrete, so it can be stored and stepped exactly. It must be deterministic, a law and not a dice roll. Its one beat must be reversible in the closed bulk, the interior of the mesh away from its growing edge, so that any arrow of time is forced to come from somewhere else. It must conserve charge, the signed sum of the tones. It must be local, each dock reading only its own sites. It must hide no state. There is one substance, not two. And it must respect the gap between grain and physics, never calling a raw tone a particle or a cluster of tones a force. Table~\ref{tab:constraints} gathers these. Breaking any of them is a cheat, not a refinement.

\begin{table}[ht]
\centering
\begin{tabular}{@{}lp{0.6\textwidth}@{}}
\toprule
constraint & what it forbids \\
\midrule
discrete & real numbers, infinite precision, a continuum in the base \\
deterministic & randomness, hidden seeds, anything but a fixed function \\
reversible & a built-in arrow of time in the closed bulk \\
charge-conserving & creating net charge from nothing \\
local & any action at a distance \\
no hidden state & memory or flags beyond the tones \\
one substance & a second kind of stuff \\
emergence gap & calling raw tones, or their first clusters, physics \\
\bottomrule
\end{tabular}
\caption{The constraints every base element must satisfy.}
\label{tab:constraints}
\end{table}

The eight, running together under these constraints, are the flow.

The whole paper depends on the emergence gap, so it gets its own statement.

\begin{principle}[Two layers]
The base is trits, three-valued units like the bits of ordinary code but holding $-1$, $0$, or $+1$ in place of $0$ or $1$. Real numbers appear only as coarse-grained averages, one or more layers up. A raw tone is never a particle, and a cluster of tones is never a force.
\end{principle}

One line in all of this cannot be tested. The base posits that the substrate \emph{feels}, that the tone is the felt state itself and not a label for it. Everything above that line can be built, stepped, and measured. Below it sits the choice to call the structure experience rather than only to describe it. That premise is the floor of the theory, and it is marked wherever it bears weight.
